\section*{Erd\H{o}s Problem 465}

\paragraph{Statement and scope.}
Let \(N(X,\delta)\) be the largest number of points in a disc of radius
\(X\) whose pairwise distances have distance at least \(\delta\)
from every integer. For fixed \(0<\delta<1/2\), must
\(N(X,\delta)=o(X)\), and even
\(N(X,\delta)\leq X^{1/2+o(1)}\)?
Both conclusions are known; Konyagin proved \(O_\delta(\sqrt X)\).
The sublinear step is unresolved here. We prove a uniform linear
bound and a sharper bound on the collinear special case.

\paragraph{A strip and fractional-coordinate bound.}
Translate the disc to the origin. Divide the vertical interval
\([-X,X]\) into at most \(\lceil8X/\delta\rceil+1\) half-open strips
of height at most \(\delta/4\).
Divide the fractional-coordinate circle, represented by \([0,1)\),
into \(\lceil4/\delta\rceil\) half-open intervals of length at most
\(\delta/4\). Assign each point \((x,y)\) the strip containing \(y\)
and the interval containing \(\{x\}\).

Two different points with the same assignment would satisfy
\[
 |\Delta y|\leq\delta/4,\qquad
 \bigl\||\Delta x|\bigr\|\leq\delta/4.
\]
The triangle inequality gives
\[
 0\leq\sqrt{\Delta x^2+\Delta y^2}-|\Delta x|
 \leq|\Delta y|.
\]
The distance-to-integers function is \(1\)-Lipschitz, so their
Euclidean distance would have distance at most \(\delta/2<\delta\)
from an integer. This contradicts the hypothesis.
Thus every assignment contains at most one point, and
\[
 N(X,\delta)\leq
 \bigl(\lceil8X/\delta\rceil+1\bigr)\lceil4/\delta\rceil
 =O_\delta(X).
 \tag{1}
\]
This improves the elementary area-packing bound of order \(X^2\).

\paragraph{The collinear case.}
For points on a line, choose signed coordinates \(t_i\) along it.
Then
\(\||P_i-P_j|\|=\|t_i-t_j\|\).
Their residues modulo one are therefore pairwise separated by
circular distance at least \(\delta\). Ordering these residues around
the circle, every gap is at least \(\delta\), and their sum is one.
There can be at most \(\lfloor1/\delta\rfloor\) points.
In particular, unbounded constructions require variation in both
coordinates.

\paragraph{Remaining gap and sources.}
The assignments in (1) impose no useful restriction between different
vertical strips. Obtaining a vanishing fraction of occupied assignments,
or the square-root bound, requires additional geometric information.
The known result is S.~V.~Konyagin,
\emph{On the distances between points on the plane},
Mathematical Notes 69 (2001), 578--581,
\url{https://www.mathnet.ru/eng/mzm691}.
Current statement and S\'ark\"ozy's earlier sublinear estimate:
\url{https://www.erdosproblems.com/465}, checked 24 September 2026.
Neither external upper-bound theorem is used to fill the gap in this attempt.

\paragraph{Completion estimate.}
The linear bound and collinear case are proved; the main sublinear improvement is missing.
\textbf{COMPLETION: 30\%}
